\section{积分的基本性质}
\label{sec:integ-basic-properties}

计算积分时我们反复做三件事：把被积函数拆成几项分别积分，用一个容易计算的函数从上方或下方比较，把积分区域切成几块分别累计。在有限情形$\sum_if(i)\mu_i$里，这三件事分别对应求和的线性、逐项比较和分组求和，都没有条件。

积分的情形多了一个麻烦：值可以是$+\infty$，而$+\infty-(+\infty)$没有意义。所以下面每条规则都要写明范围：可积函数的运算在有限实数中进行，非负函数允许积分为$+\infty$，但不允许系数为负。

\subsection{线性性}
\label{sec:integ-linearity}

\paragraph{三种情形的线性性}
线性性要回答两件事：线性组合是否仍属于可以积分的范围，积分能否逐项计算。非负函数总可以积分，但系数必须非负，否则会产生负值；可积函数允许任意实系数或复系数，代价是要先核对线性组合仍然可积。

\begin{theorem}[积分的线性性]
\label{thm:integ-linearity}
\begin{enumerate}
 \item 设$f,g:X\to[0,+\infty]$可测，$a,b\in[0,+\infty]$，约定$0\cdot(+\infty)=0$。则下式在$[0,+\infty]$中成立：
 \begin{equation}
  \int(af+bg)\,\mathrm d\mu=a\int f\,\mathrm d\mu+b\int g\,\mathrm d\mu .
  \label{eq:integ-linearity}
 \end{equation}
 \item 设$f,g$是实值可积函数，$a,b\in\mathbb R$。则$af+bg$可积，且式~\eqref{eq:integ-linearity}成立。
 \item 设$f,g$是复值可积函数，$a,b\in\mathbb C$。则$af+bg$可积，且式~\eqref{eq:integ-linearity}成立。
\end{enumerate}
\end{theorem}

\begin{proof}
第1条。由式~\eqref{eq:integ-nonnegative-additive}，$\int(af+bg)=\int af+\int bg$，所以只需证明$\int af=a\int f$。系数有限的情形第~\ref{sec:integ-nonnegative}小节已经给出，只需处理$a=+\infty$。

函数$(+\infty)\cdot f$在$\{f>0\}$上取$+\infty$，在$\{f=0\}$上取$0$。分两种情形。若$\mu(\{f>0\})>0$，这个函数的积分为$+\infty$；由引理~\ref{lem:integ-null-sets}第3条，此时$\int f>0$，右端$(+\infty)\cdot\int f$也是$+\infty$。若$\mu(\{f>0\})=0$，左端是几乎处处为零的函数的积分，等于$0$；右端是$(+\infty)\cdot0=0$。

第2条，先看可积性。可积函数只在零测集上取$\pm\infty$，把这些点上的值改成$0$，积分不变，$af+bg$处处有定义。由三角不等式与第1条，
\[
 \int\lvert af+bg\rvert\,\mathrm d\mu\leq\lvert a\rvert\int\lvert f\rvert\,\mathrm d\mu+\lvert b\rvert\int\lvert g\rvert\,\mathrm d\mu<+\infty,
\]
所以$af+bg$可积。

再看数乘。$a\geq0$时$(af)^\pm=af^\pm$；$a<0$时$(af)^+=\lvert a\rvert f^-$，$(af)^-=\lvert a\rvert f^+$。代入式~\eqref{eq:integ-signed}，两种情形都得到$\int af=a\int f$。

最后看加法。令$h=f+g$。把$h^+-h^-=f^+-f^-+g^+-g^-$中带负号的项移到另一边，得到
\[
 h^++f^-+g^-=h^-+f^++g^+ .
\]
两边都是非负可测函数之和，由第1条逐项积分，六个积分都有限，于是可以移项，得到
\[
 \int h\,\mathrm d\mu=\int f\,\mathrm d\mu+\int g\,\mathrm d\mu .
\]

第3条。第2条中的三角不等式对复数同样成立，给出可积性。加法按实部与虚部分别用第2条。对复数$a=\alpha+\mathrm i\beta$与$f=u+\mathrm iv$，有
\[
 af=(\alpha u-\beta v)+\mathrm i(\alpha v+\beta u).
\]
由第2条分别积分实部与虚部，再合并成复数的乘积：
\[
\begin{aligned}
 \int af&=\alpha\int u-\beta\int v+\mathrm i\Bigl(\alpha\int v+\beta\int u\Bigr)\\
 &=(\alpha+\mathrm i\beta)\Bigl(\int u+\mathrm i\int v\Bigr)=a\int f.
\end{aligned}
\]
\end{proof}

\paragraph{条件不能省略}
第2条要求两个函数都可积，只知道各自的积分存在是不够的。在$\mathbb R$上取$f\equiv1$、$g\equiv-1$，两者的积分分别为$+\infty$与$-\infty$，都存在。而$f+g\equiv0$的积分为$0$，右端却是$+\infty+(-\infty)$，没有意义。

第1条要求系数非负，同样是为了不让正负贡献相遇。由第2条，可积函数全体对线性组合封闭，积分是其上的线性映射，就像有限维时$f\mapsto\sum_if(i)\mu_i$是一个线性函数。这一结构在第~\ref{chap:function-normed-spaces}章作为函数空间再研究。

\subsection{单调性}
\label{sec:integ-monotonicity}

\paragraph{逐点比较}
用一个已知的函数控制被积函数，是估计积分最常用的办法。非负函数的单调性已由定义得到；对实值可积函数，比较可以借助线性性完成：$f\leq g$意味着$g-f\geq0$，非负函数的积分非负。

\begin{theorem}[积分的单调性]
\label{thm:integ-monotonicity}
设$f,g$都是非负可测函数，或都是实值可积函数。若$f\leq g$几乎处处成立，则$\int f\,\mathrm d\mu\leq\int g\,\mathrm d\mu$。若还有$\int f=\int g$且两者有限，则$f=g$几乎处处成立。
\end{theorem}

\begin{proof}
设$N$是零测集，$f\leq g$在$N$外成立。把$f,g$在$N$上都改成$0$，由引理~\ref{lem:integ-null-sets}第4条（实值情形对正负部分别使用），积分不变，而改后的函数处处满足$f\leq g$。

非负情形由定义的单调性得结论。可积情形中，$g-f$是几乎处处非负的可积函数，由线性性，
\[
 \int g\,\mathrm d\mu-\int f\,\mathrm d\mu=\int(g-f)\,\mathrm d\mu\geq0 .
\]

现在设两个积分相等且有限。非负情形下$f,g$的积分有限，也属于可积函数，所以统一按可积情形处理。这时$\int(g-f)=0$，而$g-f$几乎处处非负，由引理~\ref{lem:integ-null-sets}第3条，$g-f=0$几乎处处成立。
\end{proof}

逐点不等式放宽到几乎处处，是因为积分看不见零测集上的取值。例如在$[0,1]$上，$\mathbf 1_{\mathbb Q}\leq0$只在零测的有理点之外成立，结论$\int\mathbf 1_{\mathbb Q}\leq0$仍然正确，再结合$\mathbf 1_{\mathbb Q}\geq0$，得$\int\mathbf 1_{\mathbb Q}=0$。

\paragraph{复值积分不比较大小}
复数之间没有与运算相容的次序，所以复值积分没有单调性可言。能比较的是实值量：实部与虚部满足$\operatorname{Re}\int f=\int\operatorname{Re}f$，可以对它们分别使用单调性；模$\lvert f\rvert$是非负函数，也可以比较。下一小节就是用实部与模的比较来估计复值积分。

\subsection{积分的绝对值估计}
\label{sec:integ-absolute-estimate}

\paragraph{实值与复值情形}
积分是加权累计，取绝对值之后，正负或不同方向的贡献不再抵消，累计只会变大。这与向量的三角不等式$\lVert\sum_iv_i\rVert\leq\sum_i\lVert v_i\rVert$是同一件事：把许多向量首尾相接，总位移不超过路程之和。

\begin{theorem}[积分的绝对值估计]
\label{thm:integ-absolute-estimate}
设$f$是实值或复值可积函数，则
\begin{equation}
 \Bigl\lvert\int f\,\mathrm d\mu\Bigr\rvert\leq\int\lvert f\rvert\,\mathrm d\mu .
 \label{eq:integ-absolute-estimate}
\end{equation}
\end{theorem}

\begin{proof}
实值情形。由$-\lvert f\rvert\leq f\leq\lvert f\rvert$与单调性、线性性，
\[
 -\int\lvert f\rvert\,\mathrm d\mu\leq\int f\,\mathrm d\mu\leq\int\lvert f\rvert\,\mathrm d\mu .
\]

复值情形。记$z=\int f\,\mathrm d\mu$。若$z=0$，左端为零，不等式成立。若$z\neq0$，令$\alpha=\bar z/\lvert z\rvert$，这是模为一的复数，使$\alpha z=\lvert z\rvert$为正实数。由复线性性，
\[
 \lvert z\rvert=\alpha\int f\,\mathrm d\mu=\int\alpha f\,\mathrm d\mu .
\]
左边是实数，所以它等于右边的实部。再用复数的实部不超过它的模这一事实与单调性，并利用$\lvert\alpha\rvert=1$，得到
\[
 \lvert z\rvert=\int\operatorname{Re}(\alpha f)\,\mathrm d\mu\leq\int\lvert\alpha f\rvert\,\mathrm d\mu=\int\lvert f\rvert\,\mathrm d\mu .
\]
\end{proof}

证明中乘以$\alpha$的作用是把积分旋转到正实轴上，这样复数的模就变成了可以比较大小的实部。例~\ref{ex:integ-complex-cancel}中$\lvert\int f\rvert=0<1=\int\lvert f\rvert$，说明不等式可以严格成立：四个向量首尾相接回到原点，总位移为零，路程却是$1$。

\paragraph{误差控制}
把估计用于两个可积函数之差，就得到积分误差的界：
\begin{equation}
 \Bigl\lvert\int f\,\mathrm d\mu-\int g\,\mathrm d\mu\Bigr\rvert\leq\int\lvert f-g\rvert\,\mathrm d\mu .
 \label{eq:integ-error-bound}
\end{equation}
用$g$近似$f$时，积分结果的误差不超过两者差的模的积分。

用例~\ref{ex:integ-x-dyadic}的二进分层函数检验这个界。在$[0,1]$上用$\varphi_n$近似$f(x)=x$，逐点有$0\leq f-\varphi_n\leq2^{-n}$，所以式~\eqref{eq:integ-error-bound}的右端不超过$2^{-n}$。实际误差是
\[
 \frac12-\frac12\bigl(1-2^{-n}\bigr)=2^{-n-1},
\]
确实在这个界之内。第~\ref{sec:integ-dominated-convergence}小节将用同一个量$\int\lvert f_n-f\rvert$判断一列近似的积分是否收敛。

\subsection{积分区域的可加性}
\label{sec:integ-domain-additivity}

\paragraph{子集上的积分}
在可测子集$A$上积分，就是只累计$A$内的贡献。用指示函数把$A$外的值清零，就可以沿用全空间上的定义。

\begin{definition}[可测子集上的积分]
\label{def:integ-subset}
设$A\in\mathcal A$，$f$非负可测或可积。定义
\begin{equation}
 \int_Af\,\mathrm d\mu=\int_Xf\mathbf 1_A\,\mathrm d\mu .
 \label{eq:integ-subset}
\end{equation}
对Lebesgue测度和区间$A=[a,b]$，也写作$\int_a^bf(x)\,\mathrm dx$。
\end{definition}

这个积分也等于把$f$限制到$A$上、在$A$的可测子集族上对$\mu$所作的积分，因为两种算法在简单函数上逐层相同，再由递增逼近推广到一般函数。若$f$可积，由$\lvert f\mathbf 1_A\rvert\leq\lvert f\rvert$，$f$在每个可测子集上可积。

\paragraph{有限拆分与零测集}
设$A,B\in\mathcal A$不交。则$f\mathbf 1_{A\cup B}=f\mathbf 1_A+f\mathbf 1_B$，由线性性，
\begin{equation}
 \int_{A\cup B}f\,\mathrm d\mu=\int_Af\,\mathrm d\mu+\int_Bf\,\mathrm d\mu .
 \label{eq:integ-domain-additive}
\end{equation}
这对非负可测函数在$[0,+\infty]$中成立，对可积函数在有限数中成立；归纳可得有限个不交区域的情形。

若$N$是零测集，$f\mathbf 1_{A\cup N}$与$f\mathbf 1_A$只在$N$上可能不同，所以$\int_{A\cup N}f=\int_Af$。增删零测集不改变积分，区间端点开闭也就无关紧要：$\int_{[a,b]}f=\int_{(a,b)}f$。

\paragraph{可数拆分}
可数个不交区域$A_1,A_2,\dots$的情形不能从有限拆分直接得到，因为需要交换积分与无限求和。对非负可测函数，第~\ref{sec:integ-monotone-convergence}小节的单调收敛定理将给出
\begin{equation}
 \int_{\bigcup_nA_n}f\,\mathrm d\mu=\sum_n\int_{A_n}f\,\mathrm d\mu,
 \label{eq:integ-domain-countable}
\end{equation}
两边允许为$+\infty$；对可积函数，再由支配收敛得到同样的等式且级数绝对收敛。换句话说，对非负可测函数$f$，$A\mapsto\int_Af\,\mathrm d\mu$本身是一个测度。由非负函数生成的这类测度怎样反过来刻画，放在第~\ref{chap:measure-densities}章。
